% Section: psi'(a/5) - zeta(2), a = 1,2,3,4, by digamma differences with shifts a/5.
% Included from lambda-irrationality.tex; uses \LamFifths, \MuFifths etc. defined there.
\section{The numbers $\psi'(\frac a5)-\zeta(2)$}\label{sec:fifths}

The shifts $s=\frac a5$, $a\in\{1,2,3,4\}$, give four further numbers by the method of Section~\ref{sec:thirds}. Here the trigamma values
are not combinations of $\zeta(2)$ and a single Dirichlet $L$-value, and we work with them directly:
\[
 \psi'(s)-\zeta(2)=\sum_{n\ge0}\frac1{(n+s)^2}-\frac{\pi^2}6 .
\]
Write $c_a=\psi'(\frac a5)-\zeta(2)$. The multiplication formula gives $\sum_{a=1}^5\psi'(\frac a5)=25\zeta(2)$, so
$c_1+c_2+c_3+c_4=\frac{10\pi^2}3$, and the reflection formula $\psi'(s)+\psi'(1-s)=\pi^2/\sin^2\pi s$ gives
\[
 c_1+c_4=\frac{25+6\sqrt5}{15}\,\pi^2,\qquad c_2+c_3=\frac{25-6\sqrt5}{15}\,\pi^2 .
\]
A numerical search (PSLQ, $60$ digits) finds no rational linear relation among $1$ and any two of the $c_a$, so we treat
Theorem~\ref{thm:fifths} as four separate statements.

\subsection{Frames and kernels}
Let $s=\frac a5$, $u=5-a$, and put
\[
\begin{gathered}
 F_s(z)=\psi(z+1)-\psi(z+s),\qquad \omega=\tfrac1{25},\\
 \xi_s:=\omega F_s'(0)=-\frac1{25}\bigl(\psi'(s)-\zeta(2)\bigr).
\end{gathered}
\]
Then $k-1+s=(5k-u)/5$, and from $\psi(w+1)-\psi(w)=1/w$,
\[
\begin{gathered}
 F_s(k)-F_s(k-1)=\frac1k-\frac5{5k-u}=-\frac{u}{k(5k-u)},\\
 F_s'(k)-F_s'(k-1)=-\frac1{k^2}+\frac{25}{(5k-u)^2}.
\end{gathered}
\]
Telescoping gives the identity \eqref{eq:shift} with $\omega\Lw_{F_s}[R,Q]=M[R,Q]+\xi_sE_{R,Q}$ and the symmetric rational kernel
\[
\begin{gathered}
 M[a',b']=-\frac{u}{25(b'-a')}\sum_{k=a'+1}^{b'}\frac1{k(5k-u)}\quad(a'<b'),\\
 M[a',a']=-\frac1{25}\sum_{k=1}^{a'}\Bigl(\frac1{k^2}-\frac{25}{(5k-u)^2}\Bigr).
\end{gathered}
\]
The geometries, the members $P_J$ ($J=C$ and $J=C\setminus\{c\}$) and the admissibility condition \eqref{eq:thirds-adm},
$r\le8m\le r+9$, are those of Section~\ref{sec:thirds}.

\subsection{Content}
\begin{theorem}[content]\label{thm:content-fifths}
Let $s=\frac a5$ with $a\in\{1,2,3,4\}$ and let $r,m$ satisfy \eqref{eq:thirds-adm}. For every $J\subseteq C$ with $|C\setminus J|\le1$
there is a positive integer $d_{r,m}$, independent of $J$, with $d_{r,m}P_J\in\Z[x]$ and
\[
 \log d_{r,m}\le(\LamFifths+o(1))\,r^2,
\]
uniformly in $m$.
\end{theorem}

The denominators of $M$ divide $k(5k-u)$ and $b'-a'$. For a prime $p\ge7$ the valuation analysis of \S\ref{sec:content} and
Section~\ref{sec:thirds} applies with $\Phi(k)=k(5k-u)$. For a prime power $q=p^f$ the set $\{k\ge1:q\mid\Phi(k)\}$ consists of the
classes $0$ and $k_q$, where $5k_q=u+j_qq$ with $j_q\in\{1,2,3,4\}$. So $k_q/q$ is, up to $O(1/q)$, one of
$\kappa\in\{\frac15,\frac25,\frac35,\frac45\}$, determined by the class of $q$ modulo $5$, and by the prime number theorem in
arithmetic progressions each value of $\kappa$ occurs for primes of density $\frac14$. Proposition~\ref{prop:dual} bounds the
assignment value by an explicit function of $r$, $m$, $q$ and $k_q$, with no new dual lemma. With $\theta=q/r$, the bound for
$\theta\le\frac14$ is that of Section~\ref{sec:thirds}, with integral $\frac12+\frac1{32}+\frac{\log2}{32}$. For $\theta>\frac14$
it is piecewise linear; for $\kappa=\frac15$ the pieces are
\[
\begin{gathered}
 \tfrac{17}8\ \text{on}\ \bigl(\tfrac14,\tfrac{15}{16}\bigr],\quad
 \tfrac{13}4-\tfrac65\theta\ \text{on}\ \bigl(\tfrac{15}{16},1\bigr],\quad
 \tfrac72-\tfrac75\theta\ \text{on}\ \bigl(1,\tfrac54\bigr],\\
 \tfrac94-\tfrac25\theta\ \text{on}\ \bigl(\tfrac54,5\bigr],\quad
 \tfrac54-\tfrac15\theta\ \text{on}\ \bigl(5,\tfrac{25}4\bigr],
\end{gathered}
\]
and the integrals over $\theta>\frac14$ are
\[
 \kappa=\tfrac15:\ \tfrac{7653}{1280},\qquad
 \kappa=\tfrac25:\ \tfrac{2637}{896},\qquad
 \kappa=\tfrac35:\ \tfrac{6149}{3072},\qquad
 \kappa=\tfrac45:\ \tfrac{451}{288}.
\]
Averaging over the four classes,
\[
\begin{aligned}
 \Lambda_{\rm trop}&=\frac14\Bigl(\frac{7653}{1280}+\frac{2637}{896}+\frac{6149}{3072}+\frac{451}{288}\Bigr)
 +\frac12+\frac1{32}+\frac{\log2}{32}\\
 &=\frac{4714081}{1290240}+\frac{\log2}{32}=3.67530\ldots,
\end{aligned}
\]
the same for all four values of $a$. (The same computation gives $\frac{16433}{5120}+\frac{\log 2}{32}$ for the thirds and $\frac{12527}{3360}+\frac{\log 2}{32}$ for
the quarters.) The primes $2$, $3$ and $5$, the prime powers with $f\ge2$ and the finitely many exceptional primes contribute $o(r^2)$. In the formal proof
the piecewise linear majorants are rational, with a small additive allowance for the rounding of $m$ and of the residue classes,
and their total is $\Lambda=\frac{195975263767}{53321748480}=\LamFifths\ldots$; this is the value used below.

\subsection{Nonvanishing and decay}
For these frames
\[
 \delta_s(z)=\frac1{25}\Bigl(\frac1{z+1}-\frac1{z+s}\Bigr),
\]
and the simple pole at $z=-1$ is again isolated in its class modulo $\Z$, because $-s\notin-1+\Z$. The proof of Theorem~\ref{thm:thirds-nv}
therefore applies word for word: for every $r\ge1$, $m\ge0$ and $x\in\Q$ some $P_J$ with $|C\setminus J|\le1$ has $P_J(x)\ne0$.

The decay estimate depends on the offset $s$ only through the $o(1)$ terms (Section~\ref{sec:thirds}). Here the content is close to that of
$G\pm5\pi^2/48$, and the certificate of \S\ref{ssec:slope-cert} applies unchanged: for admissible geometries, every $J$ with $|C\setminus J|\le1$
and real $t$ with $|t|\le e^{-cr}$,
\[
\begin{gathered}
 |P_J(\xi_s+t)|\le\exp\bigl(-(D-o(1))r^2\bigr),\\
 D=\tfrac{390557}{100000}=\Dnew,\qquad c=\tfrac{18386549}{2500000}.
\end{gathered}
\]

\subsection{Proof of Theorem~\ref{thm:fifths}}
With
\[
\begin{gathered}
 \eta=D-\Lambda=3.90557-\LamFifths\ldots=\EtaFifths\ldots,\\
 \frac c\eta=\frac{7.3546196}{\EtaFifths\ldots}=\MuFifthsExact\ldots<\MuFifths ,
\end{gathered}
\]
the deduction is the one in the proof of Theorem~\ref{thm:thirds}, using Theorem~\ref{thm:content-fifths}, the nonvanishing above and the decay estimate. It gives
$\mu(\xi_s)\le c/\eta<\MuFifths$, and $\psi'(\frac a5)-\zeta(2)=-25\xi_{a/5}$.\qed
